\documentclass[sigconf, nonacm]{acmart}

\usepackage{booktabs}

%%
%% \BibTeX command to typeset BibTeX logo in the docs
\AtBeginDocument{%
  \providecommand\BibTeX{{%
    Bib\TeX}}}

%% These commands are for a PROCEEDINGS abstract or paper.
\acmConference[Explainable and Trustworthy AI - Project]{}{}{}

\begin{document}

\title{Unsupervised Concept Discovery for Medical Vision--Language Models:\\A Rigorous Characterization of MedConcept Limits to a 2D BiomedCLIP Adaptation}

\author{Marc'Antonio Lopez}
\affiliation{%
  \institution{Politecnico di Torino}
  \city{Turin}
  \country{Italy}
}

\author{Nicolò Colle}
\affiliation{%
  \institution{Politecnico di Torino}
  \city{Turin}
  \country{Italy}
}

\author{Carmine Francesco Benvenuto}
\affiliation{%
  \institution{Politecnico di Torino}
  \city{Turin}
  \country{Italy}
}

\begin{abstract}
Medical Vision--Language Models (VLMs) are accurate but opaque. MedConcept~\cite{haque2026medconcept} targets \emph{3D volumetric} abdominal CT with the Merlin~\cite{blankemeier2024merlin} foundation model---a resource-intensive setting. We \emph{adapt} its paradigm---sparse-autoencoder (SAE) decomposition of a frozen VLM's embeddings, alignment to a radiology vocabulary, LLM-as-judge evaluation---to a \emph{2D}, GPU-accessible setting on BiomedCLIP~\cite{zhang2023biomedclip}, at two scales: IU X-Ray~\cite{demner2016iuxray} (${\sim}7$k) and ROCOv2~\cite{rueckert2024rocov2} (${\sim}80$k, $10{\times}$). This deliberate accessibility trade yields compressed 2D embeddings. We first expose the slot-wise cross-seed Jaccard stability headline as \emph{floor-by-construction} (a relabelled single SAE scores ${\approx}0.0077$) and replace it with a permutation-invariant \emph{matched} metric against a subspace-conditioned null. Under it, the 2D adaptation is genuinely non-identifiable: best-match cosine $0.30$--$0.33$ at only $1.67$--$1.81{\times}$ the null, ${\approx}0\%$ matched at ${\geq}0.9$, naming $0.40$--$0.48$; neither a $10{\times}$ corpus nor a pre-projection variant lifts this, localising the constraint to the representation site---the price of decomposing a 2D compressed, contrastive-shaped global vector. SPLiCE~\cite{bhalla2024spliece}, a deterministic fixed-dictionary decomposition, is the sole judge-ready method ($81.6\%$ Aligned, MedGemma-4B vs.\ $3.3\%$, Llama-3.1-8B). We deliver a rigorous, partially-negative characterisation of an accessible 2D adaptation---consistent with MedConcept's own report of median Aligned$<0.1$ on one of its datasets.
\end{abstract}

\keywords{Concept-based explainability, sparse autoencoders, medical vision--language models, interpretability evaluation, reproducibility}

\maketitle

\section{Introduction}
Medical VLMs attain strong performance on pathology classification, segmentation, and report generation, yet their internal representations are high-dimensional, polysemic, and opaque---a serious concern in safety-critical clinical use that classical post-hoc tools (saliency, attention) rarely resolve into semantic structure.

Concept-based explainability reframes an explanation in terms of human-interpretable units. Where supervised approaches such as Concept Bottleneck Models~\cite{koh2020concept} and TCAV~\cite{kim2018tcav} require a pre-defined, labelled concept set, recent work attempts \emph{unsupervised} discovery: decomposing a frozen model's representations into sparse features, naming each feature against a clinical vocabulary, and validating the result quantitatively. MedConcept~\cite{haque2026medconcept} instantiates this as a three-stage pipeline (sparse-activation extraction, vocabulary alignment, LLM-judge evaluation) with the Aligned/Unaligned/Uncertain metrics---but targets \emph{3D volumetric} abdominal CT via the Merlin~\cite{blankemeier2024merlin} foundation model, a compute-heavy setting out of reach for course-level hardware.

We make a deliberate, accessibility-driven choice: \emph{adapt} the MedConcept paradigm to a 2D setting built on BiomedCLIP~\cite{zhang2023biomedclip} and 2D radiographs, which runs on a single consumer GPU. The trade-off is explicit---2D projected embeddings are cheaper to extract but far more compressed than Merlin's volumetric features---so we treat this as a study of \emph{where and why} the accessible 2D adaptation succeeds or falls short, attributing any shortfall to the representation rather than to a training bug. We ask whether the concepts such a pipeline discovers are \emph{reproducible} and \emph{clinically valid}, and find---rigorously, against analytical nulls---that the learned 2D sparse decomposition is essentially non-identifiable: independently seeded SAEs share no canonical features, the naming signal is weak, and many features never fire. Crucially, we first expose that the slot-wise index Jaccard used as a stability headline is \emph{floor-by-construction} (a relabelled copy of one SAE scores ${\approx}0.0077$ on it), so non-identifiability must be measured permutation-invariantly.

We probe three mitigations (a 512-d \textbf{baseline} SAE; \textbf{Path~A}, a 768-d pre-projection variant; \textbf{Path~B}, SPLiCE's fixed dictionary) plus a concept-\textbf{organisation} extension, and contribute: a \emph{2D, GPU-accessible adaptation} of MedConcept with an analysis of why it underperforms the 3D/Merlin setting; a \emph{relabeling control} showing the slot-wise Jaccard is floor-by-construction; a permutation-invariant \emph{matched-stability} metric with a subspace-conditioned null under which the concepts are non-identifiable; a $10{\times}$ scale refutation; and an honest, partially-negative evaluation.

\section{Related work}
\textbf{Concept-based explainability.} Concept Bottleneck Models~\cite{koh2020concept} restructure a network to predict an explicit concept layer before the classifier, trading minor accuracy for auditable, concept-level intervention. TCAV~\cite{kim2018tcav} tests directional sensitivity of a logit to user-defined Concept Activation Vectors. Both presuppose supervised, pre-defined concepts---precisely the limitation unsupervised discovery targets.

\textbf{Sparse autoencoders for monosemantic decomposition.} Bricken et al.~\cite{bricken2023monosemanticity} showed that overcomplete dictionary-learning autoencoders trained on transformer activations yield individual latents that activate on coherent, human-nameable features. Gao et al.~\cite{gao2024scaling} scaled the Top-K (k-sparse) autoencoder, which sets the $L_0$ sparsity exactly via top-$k$ gating and removes the L1 sparsity penalty, improving the reconstruction--sparsity frontier; we adopt this architecture.

\textbf{SAE stability and universality.} Lan et al.~\cite{lan2024universality} showed SAE feature spaces across models are similar only up to rotation (a shared subspace, not identical features), and Leask et al.~\cite{leask2025canonical} argued SAEs ``do not find canonical units of analysis.'' This weak-vs-strong universality motivates our matched-stability metric: we expect only a shared subspace, not identical features.

\textbf{Sparse concept decomposition.} SPLiCE~\cite{bhalla2024spliece} bypasses a learned dictionary and decomposes a frozen CLIP embedding directly into a sparse, non-negative combination over a fixed, interpretable concept dictionary. Because the dictionary is fixed and the method deterministic, it sidesteps the non-identifiability of learned factorisation---a property we exploit as a data-frugal baseline.

\textbf{Medical VLMs and the reference framework.} BiomedCLIP~\cite{zhang2023biomedclip} is a domain-specific contrastive VLM (ViT-B/16 + PubMedBERT, PMC-15M) whose image encoder exposes a 768-d pre-projection hidden state; its checkpoint projects this to the 512-d shared space we decompose. MedConcept~\cite{haque2026medconcept} provides the template (sparse extraction, vocabulary alignment, LLM-judge) we adapt; it operates on \emph{3D volumetric} abdominal CT with the Merlin~\cite{blankemeier2024merlin} foundation model (far more compute- and data-demanding than our 2D setting). RadLex~\cite{langlotz2006radlex} and MeSH~\cite{nlm2026mesh} the external radiology lexicons for naming.

\section{Research gaps}
From the literature and the project brief we identify five gaps that this work addresses:
\begin{enumerate}
  \item \textbf{Instability / non-identifiability:} sparse factorisations are not reproducible across seeds; no prior medical-VLM work reports seed-stability against a null.
  \item \textbf{Clinical validity:} a cosine-assigned name is cosmetic unless it tracks real pathology, yet concepts are rarely validated against ground truth.
  \item \textbf{Representation-location mismatch:} the SAE literature operates on raw hidden states, but reference pipelines fit the SAE on the already-projected CLIP space.
  \item \textbf{Small-data robustness:} learned SAEs need $10^5$--$10^6$ activations; medical corpora are 2--3 orders smaller, yielding dead features.
  \item \textbf{Non-reproducible evaluation:} single-metric claims never report a chance floor (slot-wise overlap is degenerate under per-seed permutation); metrics must be null-calibrated and permutation-invariant.
\end{enumerate}

\section{Methodology}
\subsection{Backbone, vocabulary, and split}
We use BiomedCLIP (ViT-B/16, 768-d hidden state projected by frozen $\mathrm{visual\_projection}: \mathbb{R}^{768}\!\to\!\mathbb{R}^{512}$ into shared contrastive space; images $224{\times}224$, L2-normalised). The shared 512-d space exhibits a \emph{modality gap}~\cite{liang2022mindthegap}; we estimate it as the difference of the mean image and vocabulary embeddings and subtract it before cosine naming/decomposition.

\textbf{Vocabulary (external, eval-independent).} Because naming must not be circular with the judge's ground truth, we use two \emph{external} lexicons chosen to match each corpus's domain: \textbf{RadLex}~\cite{langlotz2006radlex}, filtered to chest terms (39 anchor queries across 13 sub-domains, $+14$ NIH ChestX-ray14 seeds, ${\sim}1{,}030$ terms), for the chest X-ray domain of IU X-Ray; and \textbf{MeSH}~\cite{nlm2026mesh} (NLM; radiology branches A/C/E, ${\sim}1{,}024$ terms) for ROCOv2's broader multimodal radiology, where a chest-specific lexicon is inappropriate. Both are built by encoding terms with the text encoder and ranking by cosine to the anchors; both are free and licence-free (a deliberate alternative to the licence-gated UMLS the reference uses).

\textbf{Corpora and split.} IU X-Ray~\cite{demner2016iuxray} (7,470 frontal/lateral images across 3,852 radiographic studies) is the primary, small-scale setting; ROCOv2 (${\sim}80$k images) is a $10{\times}$-scale replication that tests the data-volume hypothesis directly. To avoid leakage we split IU X-Ray \emph{by radiographic study} (key $=\texttt{patient\_IM-study}$), so no study appears in both partitions; ROCOv2 is split by image (no shared patient structure). The IU split yields 5,955 train / 1,515 test images (3,081 / 771 studies) with verified group overlap $=0$, recomputed deterministically on every run (sorted study set, fixed seed).

\subsection{Baseline: Top-K SAE on the 512-d projected space}
We train a Top-K SAE~\cite{gao2024scaling} ($\mathrm{ReLU}(W_{\mathrm{enc}}(x-b_{\mathrm{dec}})+b_{\mathrm{enc}})$, hard top-$k$ sparsification, $b_{\mathrm{dec}}$ at the geometric median) on the 512-d image embeddings. Configuration: $D{=}2048$, $k{=}32$, 8,000 steps, $\mathrm{lr}{=}5{\times}10^{-5}$, batch size 256, five seeds $(0,42,123,456,789)$. Concept naming assigns each \emph{live} (non-dead) decoder row to its closest vocabulary term by cosine similarity, after gap correction ($W_{\mathrm{dec}}\!\leftarrow\!W_{\mathrm{dec}}-\mathrm{gap}$). Per-image pseudo-reports list the top-5 active concepts with a stated dominant concept.

\subsection{Path A: Top-K SAE on the 768-d hidden state}
Path A addresses Gap~3. Instead of $\mathrm{visual\_projection}(\cdot)$, we extract the 768-d CLS token \emph{before} the projection and train an identical Top-K SAE on these raw (un-normalised) hidden states. Because the text vocabulary lives in 512-d, we bridge the two spaces with the \emph{frozen} projection: $\hat W_{\mathrm{dec}}^{512}=W_{\mathrm{dec}}^{768}W_{\mathrm{proj}}^{\top}$, gap-correct, then name as in the baseline. Dead features are detected pre-projection and cross-checked with an activation-based mask.

\subsection{Path B: SPLiCE on a fixed dictionary}
Path B addresses Gaps~1 and~4. We decompose each gap-corrected 512-d image embedding into a sparse, non-negative combination over the fixed vocabulary dictionary with a two-stage solver (SPLiCE-style, in place of its original Lasso): Orthogonal Matching Pursuit selects $k{=}32$ atoms, then non-negative least-squares re-solves the coefficients on that support. The dictionary is never estimated from the samples, so the method is deterministic by construction and has no cross-seed instability. The output schema mirrors the SAE explanations so the same judge can score it.

\subsection{Concept-organisation extension}
Addressing Gap~5, we cluster the active vocabulary into concept families by agglomerative clustering (average linkage, cosine metric). Each cluster is annotated with the most specific shared anatomical ancestor, with two degeneracy guards: leaf-root rejection and a \emph{subtree-size canopy} that rejects ancestors whose descendant count exceeds 1\% of the ontology (e.g.\ the trivially generic ``anatomical entity'' node). The stage emits structured per-image families plus organisation metrics (silhouette, redundancy reduction, coverage); it is standalone and does not feed the judge.

\subsection{Evaluation metrics}
\textbf{Stability (matched, the headline).} For each decoder row in seed~$i$ we take the best cosine match across all rows of seed~$j$~\cite{bricken2023monosemanticity,lan2024universality}, a permutation-invariant statistic. The null is a \emph{subspace-conditioned} random-vector draw: real decoder directions concentrate in a subspace of effective rank $\mathrm{erank}{=}(\Sigma\sigma)^2/\Sigma\sigma^2{\approx}357$--$363$ (not the ambient $512$), so we draw null vectors inside each decoder's top-$\mathrm{erank}$ right-singular subspace and compare against the other decoder projected into it; the isotropic (full-$512$-d) ratio is a reported lower bound.
\textbf{Stability (slot-wise, deprecated).} Cross-seed Jaccard of top-$k$ index sets is \emph{floor-by-construction}---not invariant to per-seed feature permutation (Sec.~\ref{sec:relabel}); retained only as a deprecated identical-permutation check.
\textbf{Naming.} Mean cosine of the assigned term, plus the fraction matched at ${\geq}0.7$ and ${\geq}0.9$ (random directions cannot reach $0.9$).
\textbf{Faithfulness.} Point-biserial correlation of feature activations with in-distribution IU X-Ray MeSH/Problems labels, with a per-feature shuffle null. \textbf{LLM judge.} Each concept is scored Aligned/Unaligned/Uncertain against the reference report by a prompted LLM (MedGemma-4B~\cite{sellergren2025medgemma} and, for cross-LLM comparison, Llama-3.1-8B); we report the raw \%Aligned as a measure of \emph{local} plausibility.

\section{Experiments and analysis}

\subsection{The 2D baseline is a non-identifiable failure case (Gap 1)}
Table~\ref{tab:baseline} reports the baseline on both corpora. Reconstruction is faithful (cosine $0.968$--$0.991$) yet cross-seed feature identity is absent: matched cosine $0.30$--$0.33$ at only $1.67$--$1.81{\times}$ the null, ${\leq}0.3\%$ matched at ${\geq}0.9$, naming capped at $0.40$--$0.48$ ($0\%$ above $0.7$). This is genuine non-identifiability (init distinct, decoders full-rank, no collapse); the deprecated slot-wise Jaccard lands at its ${\approx}0.0077$ floor (Sec.~\ref{sec:relabel}).

\begin{table}[t]
\caption{Baseline Top-K SAE ($D{=}2048,k{=}32$); IU X-Ray (RadLex) vs.\ ROCOv2 (MeSH, $10{\times}$). ``Sub./iso.\ null'' = observed best-match over the subspace-conditioned/isotropic null.}
\label{tab:baseline}
\small
\begin{tabular}{lcc}
\toprule
 & IU X-Ray (1$\times$) & ROCOv2 (10$\times$) \\
\midrule
Reconstruction cosine & 0.991 & 0.968 \\
Dead features (full test) & 15.0--17.8\% & 0.5--0.8\% \\
Matched cosine (obs) & 0.299 & 0.327 \\
Decoder erank & 363 & 357 \\
Subspace null ratio (honest) & 1.67$\times$ & 1.81$\times$ \\
Isotropic null ratio (lower bnd) & 1.98$\times$ & 2.16$\times$ \\
Frac.\ matched $\geq0.9$ & 0.0\% & 0.3\% \\
Naming cosine (live, mean) & 0.40 & 0.48 \\
Slot-wise Jaccard (deprecated) & 0.0077 & 0.0077 \\
\bottomrule
\end{tabular}
\end{table}

\subsection{Relabeling control: the slot-wise metric is floor-by-construction}
\label{sec:relabel}
The slot-wise Jaccard cannot distinguish ``two different SAEs'' from ``one SAE whose columns were renamed.'' Permuting one trained ROCOv2 SAE's rows by a random bijection~$\pi$ (function-preserving; reconstruction unchanged to $10^{-4}$) makes slot-wise Jaccard collapse to $0.0077$ (indistinguishable from real cross-seed) while the matched metric correctly reports identity ($1.0$, frac${\geq}0.9{=}1.0$). The slot-wise signal is pure permutation noise---a binary identical-permutation check, not a stability metric. All later stability claims use the matched metric.

\subsection{Neither representation depth nor data scale restores identifiability}
Table~\ref{tab:patha} contrasts Path~A (768-d pre-projection hidden state) with the baseline. Operating on the 768-d hidden state markedly reduces dead features ($1.7\%$ vs.\ $16\%$) and improves naming cosine ($0.47$ vs.\ $0.42$). Yet the permutation-invariant matched metric is unchanged in its verdict: both share a decoder subspace above the null (isotropic $1.97{\times}$/$2.63{\times}$, collapsing to ${\approx}1.3$--$1.5{\times}$ under the subspace-conditioned null), and almost no features match strongly ($0\%$ at ${\geq}0.9$)---\emph{weak universality}~\cite{lan2024universality,leask2025canonical}: a shared subspace exists, but not canonical, reproducible features. Path~A still operates on the \emph{pooled CLS token}, so the representation-site mismatch is only partially addressed.

We then test data scale (Gap~4) directly. Training the baseline on ROCOv2 (${\sim}80$k, $10{\times}$ over IU) improves training health (dead features $16\%{\to}0.6\%$) but leaves identifiability flat (Table~\ref{tab:baseline}): matched cosine $0.299{\to}0.327$, subspace ratio $1.67{\times}{\to}1.81{\times}$, frac${\geq}0.9$ $0.0\%{\to}0.3\%$. A $10{\times}$ corpus lifts none of these above the non-identifiability regime.

\begin{table}[t]
\caption{Baseline (512-d) vs.\ Path~A (768-d), parity config.; ``obs/null'' over the isotropic null.}
\label{tab:patha}
\small
\begin{tabular}{lcc}
\toprule
 & Baseline 512-d & Path A 768-d \\
\midrule
Slot-wise Jaccard & 0.0084 & 0.0083 \\
Analytical floor & 0.0079 & 0.0079 \\
Dead features (full test) & 16.0\% & 1.7\% \\
Naming cosine (live) & 0.42 & 0.47 \\
Matched cosine (obs) & 0.299 & 0.325 \\
Matched cosine (null) & 0.151 & 0.124 \\
obs/null & 1.97$\times$ & 2.63$\times$ \\
Frac.\ matched $\geq0.9$ & $\approx$0\% & $\approx$0\% \\
\bottomrule
\end{tabular}
\end{table}

A dictionary/$k$ sweep on Path~A (Table~\ref{tab:ablation}) confirms the pattern: matched obs/null falls from $2.78{\times}$ ($D{=}1024$) through $2.63{\times}$ (default) to $2.00{\times}$ ($D{=}4096$) as overcompleteness grows---no setting escapes weak universality.

\begin{table}[t]
\caption{Path~A dictionary/$k$ sweep (768-d, 8k steps).}
\label{tab:ablation}
\small
\begin{tabular}{lcccc}
\toprule
Preset & $D$ & $k$ & Dead & Matched obs/null \\
\midrule
Conservative & 1024 & 16 & 41.7\% & 2.78$\times$ \\
Default & 2048 & 32 & 12.9\% & 2.63$\times$ \\
Aggressive & 4096 & 64 & 6.6\% & 2.00$\times$ \\
\bottomrule
\end{tabular}
\end{table}

\subsection{Why the 2D adaptation underperforms the 3D MedConcept setting}
\label{sec:hypotheses}
This negative identifiability result is specific to the accessible 2D adaptation, not to the MedConcept paradigm: MedConcept decomposes Merlin's high-dimensional \emph{3D volumetric} features~\cite{blankemeier2024merlin}, whereas we decompose BiomedCLIP's 2D \emph{projected} global vector. We attribute the gap to three structural factors: \emph{(i)} Merlin is a larger 3D CT foundation model than the lighter 2D BiomedCLIP. This difference leads to a post-projection 512-d contrastive vector that has effective rank $\mathrm{erank}{\approx}357$--$363$, so a $2048$-way dictionary is over-complete and the factorisation non-identifiable, whereas Merlin's volumetric embedding is wider and higher-rank; \emph{(ii)} the SAE literature targets raw residual streams, not a pooled, L2-normalised, contrastive-shaped summary---Path~A (pre-projection CLS) only partially relieves this; lastly, \emph{(iii)} learned SAEs need $10^5$--$10^6$ activations, and even our $10{\times}$ corpus stays 1--2 orders below. The binding constraint is thus the \emph{2D compressed representation} we adopted for GPU accessibility; the experiment our diagnosis predicts would restore identifiability is an SAE on the BiomedCLIP \emph{patch-token residual stream} (mid layer, not pooled)~\cite{bricken2023monosemanticity}.

\subsection{Path B is deterministic and judge-ready}
\label{sec:pathb}
SPLiCE decomposes every test image with a fixed, randomness-free solver (OMP selects the $k{=}32$-atom support, NNLS re-solves coefficients), so the result is deterministic by construction with no cross-seed instability, and its SAE-compatible schema lets the same judge read SPLiCE concepts. On IU X-Ray it covers all 1,515 test images in 9.5\,s using 997/1,031 RadLex terms ($96.7\%$, $14.18$ concepts/image); on ROCOv2 (${\sim}80$k) it covers all 15,958 images in 101\,s using 1,024/1,024 MeSH terms ($100\%$, $8.56$ concepts/image). Gap correction raises the maximum atom--image cosine sharply on both corpora (IU $0.49{\to}0.85$, ROCOv2 $0.54{\to}0.86$), confirming the gap, not the solver, was binding; frequent terms are mixed (a property of the supplied dictionaries, not of SPLiCE).

\subsection{Concept organisation is dense but weakly separated}
Applied to all three IU sources (Table~\ref{tab:org}), SPLiCE groups 997 concepts into 32 families, cutting concepts-per-image $14.18{\to}8.09$ ($1.75{\times}$) but with weak silhouette ($0.020$); the noisy IU vocabulary makes every cluster inherit a noisy ancestor. The rebuilt SAE baseline exposes 80 active concepts (9 families, silhouette $0.094$, no empty images) that inherit the baseline's non-identifiability, and Path~A's hidden source stays degenerate (14 concepts, 1{,}260/1,515 empty). ROCOv2 SPLiCE replicates non-degenerately (1{,}024 concepts, 32 families, silhouette $0.066$): density scales with corpus size, separation stays weak.

\begin{table}[t]
\caption{Concept organisation across sources.}
\label{tab:org}
\small
\begin{tabular}{lcccc}
\toprule
Source & Active & Clusters & Silhouette & Empty img.\ (/1515) \\
\midrule
SPLiCE & 997 & 32 & 0.020 & 0 \\
SAE baseline & 80 & 9 & 0.094 & 0 \\
SAE hidden & 14 & 4 & 0.284 & 1260 \\
\bottomrule
\end{tabular}
\end{table}

\subsection{A small upper tail of features is genuinely faithful}
Despite global instability, a faithfulness ablation ($D{=}2048$ baseline, 5 seeds) finds $17.8\%\pm0.9$ of live features exceed the shuffle-null $95$th percentile in point-biserial correlation with IU X-Ray MeSH/Problems labels ($1.1\%$ at $|r|>0.30$; strongest $|r|{=}0.459\pm0.057$ on a \emph{mass} feature)---a real but small upper tail atop an otherwise degenerate dictionary: concepts are simultaneously \emph{unstable cross-seed} (zero cluster across seeds at ${\geq}0.90$, not above a seed-tag shuffle-null, $p{=}0.17$) and \emph{partially faithful} (Gap~2).

\subsection{The judge finds high but model-dependent local plausibility}
\label{sec:judge}
Running the LLM-judge to convergence ($\mathrm{temperature}{=}0$, zero parse errors) with two backbones---MedGemma-4B~\cite{sellergren2025medgemma} and Llama-3.1-8B---SPLiCE on IU X-Ray scores $81.6\%$ Aligned under MedGemma ($1{,}230/1{,}508$) but only $3.3\%$ under Llama; the degenerate Path~A hidden source scores $76.3\%$ vs.\ $0.5\%$; the SAE baseline on ROCOv2 $88.3\%$ ($14{,}095/15{,}957$) vs.\ $23.1\%$ (Llama leans \emph{Uncertain}). This \emph{local} plausibility is orthogonal to the global identifiability verdict and is model-dependent. Low alignment is a property of the paradigm itself: MedConcept reports median Aligned$<0.1$ on AbdomenAtlas~\cite{haque2026medconcept}---so our scores are consistent with the harshness of report-conditioned concept verification (given the different backbone and modality), not evidence the 2D setting is uniquely broken.

\section{Conclusions}
We adapted MedConcept to a 2D, GPU-accessible setting (BiomedCLIP + 2D radiographs). Because the slot-wise stability headline is floor-by-construction (a relabelled SAE scores ${\approx}0.0077$), we measure non-identifiability permutation-invariantly: under a matched metric with a subspace-conditioned null, the 2D SAE concepts are non-identifiable (${\approx}0\%$ at ${\geq}0.9$, naming $0.40$--$0.48$), and neither Path~A nor a $10{\times}$ corpus lifts it---a structural limit of decomposing a 2D compressed, contrastive-shaped global vector, not a training bug (Sec.~\ref{sec:hypotheses}). SPLiCE on a fixed dictionary is the sole deterministic, judge-ready alternative, and a small (${\sim}18\%$) upper tail of features is genuinely faithful.

\textbf{Limitations and future work.} (i)~The LLM-judge returns \emph{model-dependent} verdicts, orthogonal to identifiability and not directly comparable to MedConcept's (different backbone/modality/judge); random-k null explanations score comparably ($79.6\%$ MedGemma, $2.9\%$ Llama), indicating weak discriminability. (ii)~The predicted mitigation is an SAE on the BiomedCLIP \emph{patch-token residual stream} (mid layer, not pooled CLS)~\cite{bricken2023monosemanticity}, testing whether the 2D limit is intrinsic to the representation site. (iii)~Our $10{\times}$ scale test (ROCOv2) also broadens the \emph{domain} (multimodal, not just chest), confounding scale with breadth; the clean test is PadChest (${\sim}160$k chest X-rays, same IU X-Ray domain), begun but unfinished---to test whether scale alone restores reproducibility.

%%
%% The next two lines define the bibliography style to be used, and
%% the bibliography file.
\bibliographystyle{ACM-Reference-Format}
\bibliography{biblio}


\end{document}
\endinput
%%
